%% kapa: exam=1 round=31 year=2020 period=1 form=A subject=civil_law subject_ko=민법 seq=16 number=16
%% source: Q-Net 감정평가사 기출문제 (cst00309) artlSeq=5207219 file=extracted/1008/5207219/A형/A형/2020년도 제31회 감정평가사 1차 시험_1교시_A형.pdf
%% conversion: kordoc 4.24.51 -> tools/gen1.py (auto); review=proofread-sonnet; flags=none
\begin{question}{16}
형성권에 관한 설명으로 옳은 것을 모두 고른 것은? (다툼이 있으면 판례에 따름)
\begin{framedtext}
ㄱ. 형성권의 행사는 상대방에 대한 일방적 의사표시로 한다. \par
ㄴ. 다른 사정이 없으면, 형성권의 행사에 조건 또는 기한을 붙이지 못한다. \par
ㄷ. 다른 사정이 없으면, 형성권은 그 일부를 행사할 수 있다. \par
ㄹ. 다른 사정이 없으면, 형성권은 제척기간의 적용을 받는다. \par
\end{framedtext}
\begin{choices}
  \item ㄱ, ㄴ, ㄷ
  \item ㄱ, ㄴ, ㄹ
  \item ㄱ, ㄷ, ㄹ
  \item ㄴ, ㄷ, ㄹ
  \item ㄱ, ㄴ, ㄷ, ㄹ
\end{choices}
\end{question}
\answer{2}
